\section{Future Work}\label{sec:future-work}

The most immediate next step is a more reusable access format for the corpus --- one that makes it possible to work with the data without re-running the entire pipeline, while remaining consistent with the legal and ethical constraints discussed in Section~\ref{sec:limitations}.

A second direction is linguistic enrichment. Adding annotation would make the corpus more useful for further work on non-standard Slovene, code-switching, and interaction in short-form social-media text. It would also open the way to content-oriented analyses that go beyond the descriptive focus of the present paper, including topic analysis, studies of register, lexical choice, and discourse patterns in annotated texts. Given the high share of replies in the corpus, thread-level analysis would be especially relevant for examining how interaction unfolds conversationally rather than only at the level of isolated posts.

A third direction is comparative and longitudinal research. The present corpus could be extended forward in time in order to follow later developments in the Slovene Bluesky space, including possible shifts in activity, participation, and language use. It would also be valuable to compare this corpus with other Slovene online communities and platforms, such as Twitter/X, Mastodon, or Threads, in order to examine how platform architecture, migration context, and local uptake shape interactional patterns, linking practices, and public conversation. More broadly, the same collection approach could be adapted to other small-language communities on Bluesky, supporting corpus construction for smaller language communities on emerging social-media platforms.
